% ---------------------------------------------------------------------------
% Tweet Virality Simulator (TVS) paper draft.
%
% This compiles with a plain LaTeX install (pdflatex tvs_paper.tex).
% For an ICWSM / AAAI submission, replace the documentclass and preamble with
% the official aaai style files (aaai24.sty etc.) and move the body across.
%
% Places that need your real numbers are marked with \fill{...} so they show
% up clearly. Search for "fill" before you submit.
% ---------------------------------------------------------------------------
\documentclass[11pt]{article}

\usepackage[margin=1in]{geometry}
\usepackage{times}
\usepackage{graphicx}
\usepackage{booktabs}
\usepackage{amsmath}
\usepackage{url}
\usepackage{authblk}

% A visible placeholder so unfinished numbers are easy to spot.
\newcommand{\fill}[1]{\textbf{[#1]}}

\title{Knowing Before You Post: A Self Improving Simulator that Estimates How Far a Tweet Will Travel}

\author[1]{First Author}
\author[1]{Second Author}
\affil[1]{Your Affiliation, City, Country}
\affil[ ]{\texttt{\{first, second\}@youruniversity.edu}}

\date{}

\begin{document}
\maketitle

\begin{abstract}
\noindent
Most work on tweet popularity waits until a post is already spreading and then
forecasts how big the wave will get. That is useful for trend tracking, but it
does not help a person who is staring at a draft and wondering whether it is
worth posting. We study the harder question: given only the text of a tweet
that has not been posted yet, can we say how far it is likely to spread for that
specific author? We build a simulator that turns a draft into a small synthetic
crowd, lets that crowd react and pass the message along through both the
follower graph and an algorithmic feed, and reads off how far the message
travels across many random runs. The score is trained not on raw likes, which
mostly reflect how famous the account is, but on how much a tweet beats the
author's own typical performance. A learning loop logs every prediction,
collects the real result later, refits the simulator, and only swaps in a new
version when it clearly beats the current one on held out data from a later time
period. The stack also fits per niche appeal weights, per community engagement
offsets, optional author specific crowd priors, a learned reaction model when
impression level data is available, ensemble uncertainty bands, causal A/B
experiments over profile versions, and daily drift monitoring with automated
retrain triggers. We trained and evaluated the system on \fill{100{,}000} tweets
from \fill{N} public accounts. The calibrated simulator reaches a rank
correlation of \fill{X.XX} between its score and real relative spread, compared
with \fill{Y.YY} for the uncalibrated baseline, while scoring a draft in about
\fill{Z} seconds. We release the simulation engine as open source and describe
the full loop from draft to outcome to retraining.
\end{abstract}

\section{Introduction}

Every day people decide whether a post is good enough to share. Brands, writers,
campaigners, and ordinary users all make a quiet bet each time they hit send.
The question they are really asking is simple. Will this travel, or will it sit
there with a handful of likes?

Research on information spread has mostly answered a related but different
question. Given a post that is already collecting reshares, how large will the
final cascade be? Methods built on point processes and on learned features do
this well once a post has a head start \cite{zhao2015seismic,cheng2014cascades}.
The catch is that they need the early part of the cascade as input. They cannot
score a blank draft, because there is no resharing history to read from yet.

A second line of work predicts popularity straight from content and account
features using machine learning \cite{rameez2022viralbert}. These models can run
before posting, but they tend to learn account size. A tweet from a star with
ten million followers will be predicted to do well almost regardless of what it
says. That is correct on average and useless for the person trying to improve a
draft, because they cannot change their follower count between drafts.

A newer group of papers simulates the crowd itself, often with large language
model agents, and reads spread off the simulation
\cite{liu2025popsim,park2023agents}. This is closer to what we want, but the
language model agents are slow and the systems are usually fit once on a fixed
benchmark rather than kept honest against a stream of real outcomes.

We take a middle path. Our simulator uses light, fast agents rather than
language model agents, so it can run many random trials per draft inside a normal
web request. It is built for the moment before posting, so it never needs cascade
history. And it is trained on a target that strips out account size, so the score
reflects the message and not the megaphone.

Our contributions are the following. First, a simulator that estimates spread for
an unposted tweet by running a synthetic cascade across a follower graph and an
algorithmic feed. Second, a training target we call relative amplification, which
measures how far a tweet beats its own author's normal results, so the same model
is fair to a small account and a large one. Third, a closed learning loop that
logs predictions, gathers outcomes, refits the simulator with a black box
optimizer, and promotes a new version only when it wins on data from a later time
window. Fourth, a full continual learning stack---niche appeal fitting,
community offsets, author audience profiles, a learned reaction model, ensemble
uncertainty, causal experiments, and drift monitoring---wired end to end in a
deployed service. Fifth, an evaluation on \fill{100{,}000} tweets with ablations
over every major component.

\section{Related Work}

\paragraph{Cascade forecasting.}
A long line of work predicts the eventual size of a cascade from its early
shape. SEISMIC treats resharing as a self exciting process and gives strong real
time size estimates \cite{zhao2015seismic}. Cheng and colleagues frame growth as
a classification problem and show that early breadth predicts later size
\cite{cheng2014cascades}. These methods are accurate, but they read from an
observed cascade, so they answer a question that only exists after posting.

\paragraph{Content based popularity.}
Other work predicts engagement from the post itself and from author features,
often with transformer text encoders \cite{rameez2022viralbert}. These run
before posting. Our concern with them is the target. When the label is raw reach
or raw likes, the strongest signal is follower count, and the model mostly
ranks accounts rather than messages.

\paragraph{Agent based and simulation models.}
Simulating the crowd has a long history in this field. Early work showed that
limited attention plus network structure alone can explain why a few messages go
far while most die \cite{weng2012competition}. Others modeled the decision to
reshare from feed position and author standing \cite{pezzoni2013retweet}. Recent
systems put language model agents in a sandbox and read spread or popularity off
the simulation \cite{liu2025popsim,park2023agents}. We share the simulation idea
but keep the agents simple so the model is fast and can be retrained often
against live results.

\paragraph{Relative measures of reach.}
Our training target borrows its spirit from audits of algorithmic amplification,
which compare what an algorithm shows against a neutral baseline rather than
looking at raw totals \cite{huszar2022amplification}. We move that idea down to
the level of a single author. Instead of comparing feeds, we compare a tweet
against the author's own median tweet.

\section{The Simulator}

The engine turns one tweet into a distribution over outcomes. Figure~\ref{fig:pipeline}
summarizes the flow. The subsections below walk through each stage in the order
the code runs them.

\subsection{Tweet DNA}
Before any crowd is built, we read the draft into a fixed feature vector we call
\textbf{Tweet DNA}. This is not a black box embedding. It is a small, named set
of signals that a human can inspect.

\paragraph{Semantic pull (six scores, each from 0 to 1).}
Hook strength measures whether the opening line stops the scroll. Emotional
intensity captures anger, joy, fear, and similar tone. Novelty flags surprising
or contrarian claims. Controversy marks debate bait and polarizing framing.
Relatability picks up universal experiences and direct address. Clarity reflects
how easy the sentence is to parse on a fast feed.

\paragraph{Structural flags.}
We also record length, hashtag count, mention count, whether a link is present,
whether media is attached, and whether the text is a question. These matter
because links and hashtag stuffing often suppress reach on the platform we model.

\paragraph{Topic vector.}
A 24 dimensional topic vector situates the tweet in interest space so the seed
audience can be biased toward people who care about the subject. The production
encoder is a lightweight hashing trick for speed; we also support a frozen
sentence transformer (e.g.\ MiniLM) for stronger semantic alignment when latency
budget allows. Both paths feed the same cascade interface.

\paragraph{Niche appeal.}
Beyond global DNA scores, each tweet is scored for appeal within a fixed set of
\fill{K} topical niches (technology, politics, sports, and similar anchors).
Soft membership weights combine with per niche appeal coefficients learned from
data (Section~\ref{sec:niche-appeal}). This lets the same wording land differently
depending on which communities are likely to see it.

The default DNA scorer is a fast rule based module. The same slot accepts a
language model or an external compat endpoint with no change to the simulator
core. Reaction probabilities and appeal in the cascade read directly from DNA,
so calibration can learn which dimensions actually predict spread for our data.

\subsection{Building a crowd}
For each draft we generate a synthetic audience of about one thousand personas.
Each persona carries:
\begin{itemize}
  \item an interest vector (clustered into \fill{12} communities),
  \item a latent popularity score from a heavy tailed distribution (a few hubs, many
  small accounts),
  \item personal rates for liking, replying, and resharing, drawn from learned
  Beta priors,
  \item optional per community offsets that shift those rates up or down after
  calibration.
\end{itemize}
When a known author connects through OAuth, the service builds and stores an
\textbf{audience profile} for that author: Beta birth priors for like, reply,
and share rates, engagement scale, popularity tail shape, and optional community
offsets. At predict time we resolve author profile $\rightarrow$ generic profile
$\rightarrow$ global champion and patch the crowd's birth parameters so the
synthetic followers match that author's real follower makeup.

\subsection{Building the follower graph}
We connect personas with a directed follow graph. Two forces shape who follows
whom.

\paragraph{Homophily.}
People are more likely to follow others with similar interest vectors. Similarity
is cosine distance in topic space, scaled by a learned homophily parameter. This
produces echo chambers and topical clusters without hand drawing a network.

\paragraph{Preferential attachment.}
Popular personas attract more incoming links. We combine homophily with a
popularity term so hubs exist but still mostly connect within their topical
neighborhood.

Each persona follows a random number of accounts (Poisson around a learned mean).
The graph is stored as, for each node, the list of its \emph{followers} (who sees
its posts when it shares). That list drives the in network cascade channel.

\subsection{Running the cascade (Monte Carlo)}
A single \emph{run} is one stochastic replay of the post's life. We repeat about
one hundred and twenty runs per draft and aggregate.

\paragraph{Seed.}
The author's in network followers see the tweet first. Who is seeded depends on
topic overlap with the draft's DNA vector and on persona popularity.

\paragraph{Reaction draw.}
Each exposed persona independently draws like, reply, retweet, and quote actions.
Probabilities combine four things: the persona's base rates, the tweet's DNA and
niche weighted appeal, calibrated scale parameters, and---when trained---
outputs from a small multilayer perceptron (MLP) reaction model. The MLP takes
persona traits, DNA features, and niche membership and predicts like, reply,
retweet, and quote logits; it is fit on impression level outcomes from
OAuth authorized posts and loaded at inference whenever a champion reaction
artifact exists. When no trained model is available, we fall back to a
transparent weighted formula inspired by the platform's public ranking weights,
so the cascade always runs.

\paragraph{In network spread.}
Anyone who retweets or quotes exposes the tweet to their followers who have not
seen it yet. This is the social graph channel.

\paragraph{Algorithmic spread.}
When round engagement crosses a learned threshold, the tweet is injected to
people outside the follower graph, standing in for the For You feed. A pool size
grows with engagement and exploration adds random reach. We track what share of
total spread came from the graph versus the algorithm.

\paragraph{Stop.}
A run ends when no new personas are exposed or the round cap is hit. From all
runs we collect reach, branching factor $R$, viral hit rate $p_{\text{viral}}$,
and the typical reach curve over rounds.

\subsection{Scoring}
The final score from zero to one hundred blends four signals on a log scale, so
that strong tweets spread out across the range instead of all sticking at the
top. The weights are amplification past the seed audience at \(0.30\), the
fraction of the crowd reached at \(0.25\), the branching factor at \(0.25\), and
the chance of a viral run at \(0.20\). A short verdict in words accompanies the
number.

\subsection{Uncertainty (optional ensemble)}
The API can return more than a point score. When ensemble mode is on, we report
the 10th and 90th percentile of scores across simulation seeds, a confidence
label (high, medium, or low), and a short note on why the spread is wide (Monte
Carlo noise, ambiguous DNA, or disagreement between profile versions). This
costs extra compute but no extra training data.

% Placeholder for a pipeline figure — add a PDF or PNG when you have one.
\begin{figure}[t]
\centering
\fbox{\parbox{0.92\linewidth}{\centering\vspace{2.2em}
\textbf{[Pipeline figure:]} Tweet $\rightarrow$ DNA $\rightarrow$ Audience +
Network $\rightarrow$ Monte Carlo cascade $\rightarrow$ Score\\
{\small (add \texttt{paper/figures/pipeline.pdf} and uncomment
\texttt{\textbackslash includegraphics} before submission)}\vspace{2.2em}}}
\caption{End to end flow. DNA features feed both appeal and topic biased seeding;
the follower graph and algorithmic channel carry the cascade; many runs produce a
score distribution.}
\label{fig:pipeline}
\end{figure}

\section{What We Train On}

\subsection{Relative amplification}
We deliberately avoid training on raw reach. Raw reach teaches a model that big
accounts do well, which we already know. Instead we measure how far a tweet beats
its own author's normal result. For an author with a history of posts, let the
baseline be the median result of their recent tweets. The label for a new tweet
is

\begin{equation}
\text{relative amplification} = \frac{\text{observed result}}{\max(\text{author baseline}, 1)}.
\end{equation}

A value near one means a typical tweet for that person. A value of five means the
tweet did five times better than they usually do. Because the baseline is the
author's own history, the same number means the same thing for a small account
and a large one.

We compute the observed result in one of two ways. When we have author authorized
access we use impressions, which is the cleanest signal. When we only have public
numbers we use a weighted sum of engagement in which replies count most, then
quotes, then retweets, then likes, echoing how the platform appears to value
deeper interaction. We sort the resulting amplification into a few tiers for
reporting.

\subsection{Why this matters for fairness across accounts}
This choice changes what the model can learn. A celebrity post with fifty
thousand likes might be a weak result for that celebrity, while the same number
would be the post of the year for a small account. By dividing through the
author's own baseline, the model is pushed to read the message rather than the
audience.

\section{The Learning Loop}

The system improves itself over time through a pipeline that runs in production.
Table~\ref{tab:components} maps each major piece to what it does in plain terms.

\begin{table}[t]
\centering
\caption{Major components of the training and serving stack.}
\label{tab:components}
\small
\begin{tabular}{p{0.22\linewidth}p{0.70\linewidth}}
\toprule
Component & Role \\
\midrule
Outcome snapshots &
Poll tweets at 30m, 1h, 6h, 24h, and 72h so we can train on early velocity
before the final label is ready. \\
Relative amplification label &
Train on how far a tweet beat its author's baseline, not raw likes. \\
Expanded profile search &
Fit global cascade knobs plus crowd birth priors (Beta shapes, engagement
means). \\
Community offsets &
Per community tweaks to like, reply, and share rates (CMA-ES + $L_2$);
persisted on champion artifact and merged at every predict request. \\
CMA-ES + common random numbers &
Search $\sim$\fill{55} parameters with an adaptive black box optimizer; fixed
seeds across candidates so noise does not fake improvements. \\
Champion challenger gate &
Promote a new profile only if it wins on a recent time holdout across rank
correlation, tail correlation, calibration, and saturation. \\
Author audience profiles &
Store fitted crowd priors per connected author; patch the simulator at predict
time. \\
Learned reaction model &
Small MLP on impression level OAuth outcomes; champion checkpoint loaded at
inference, heuristic fallback otherwise. \\
Ensemble uncertainty &
Optional p10/p90 score band and confidence from extra simulation runs. \\
A/B experiment layer &
Route a fraction of traffic to profile arms and compare outcomes causally. \\
Drift monitoring &
Daily checks on rank correlation, saturation, label coverage, and score
monotonicity; writes structured events and triggers retrain when thresholds
are breached. \\
Niche appeal weights &
Ridge regression on soft niche membership $\times$ DNA features, fit after
CMA-ES and stored in the champion profile payload. \\
\bottomrule
\end{tabular}
\end{table}

Table~\ref{tab:phases} maps the same stack to the staged build we used in
engineering. Every row is part of the submitted system.

\begin{table}[t]
\centering
\caption{Continual learning phases (all included in the deployed service).}
\label{tab:phases}
\small
\begin{tabular}{clp{0.58\linewidth}}
\toprule
Phase & Name & What it does \\
\midrule
1--3 & Ingest \& labels &
Log every predict; snapshot velocity at 30m/1h/6h/24h; mature at 72h;
label with relative amplification. \\
4 & Community offsets &
Per community shifts to like/reply/share rates; CMA-ES search with strong
$L_2$ penalty; stored on champion artifact and reloaded at serve time. \\
5--7 & CMA-ES \& gate &
$\sim$\fill{55} global + birth + Beta parameters; common random numbers;
champion--challenger promotion on time holdout with tail metrics. \\
8 & Embeddings &
Hashing trick default; optional sentence transformer encoder for DNA topic
vector (same cascade API). \\
9 & Reaction MLP &
Train on OAuth impression outcomes; champion artifact replaces heuristic
reaction draw when available. \\
10 & Early velocity &
High frequency snapshots feed interim labels before 72h maturation. \\
11 & Author audience &
Per author crowd priors; auto built after OAuth connect; resolved at predict. \\
12 & Ensemble &
Optional p10/p90 bands, confidence, and spread reason via extra MC runs. \\
13 & Causal A/B &
Experiment arms over profile or prompt variants; outcomes linked per arm. \\
14 & Monitoring &
Daily drift job; saturation and rank corr checks; retrain cron weekly. \\
\bottomrule
\end{tabular}
\end{table}

\paragraph{Log.}
Every time the service scores a tweet, it stores the draft, the DNA features, and
the predicted score.

\paragraph{Collect.}
Later, the service gathers real metrics on a schedule (snapshots first, final
label at 72h). Results arrive through author authorized access (impressions) or
public engagement proxies.

\paragraph{Refit.}
We fit the simulator's parameters to the paired data in two stages. First,
CMA-ES \cite{hansen2016cma} searches global cascade knobs, crowd birth priors
(Beta shapes, engagement log mean), and community offset vectors. Every
candidate profile is scored on the training slice with the \emph{same} Monte
Carlo seeds (common random numbers), so the optimizer sees real differences and
not lottery luck. The objective rewards overall rank match, tail rank match, top
decile recall, and penalizes scores stuck at 0 or 100. Second, we run niche
appeal fitting (Section~\ref{sec:niche-appeal}) on the same training slice and
merge the resulting per niche weights into the challenger profile before
evaluation.

\subsection{Niche appeal fitting}
\label{sec:niche-appeal}
Global DNA appeal treats every audience the same. In practice, a finance joke and
a sports meme need different coefficients. After CMA-ES, we hold the cascade
parameters fixed and solve a ridge regression over niche membership. For each
training tweet we compute soft membership $M_{ik}$ over $K$ fixed niche
anchors and DNA feature vector $F_i$. We stack rows $M_{ik} \cdot F_i$ into a
design matrix and regress toward rank standardized relative amplification, with
an $L_2$ prior pulling weights back to the global appeal vector. The fitted
$K \times D$ weight matrix is written into \texttt{niche\_appeal\_weights} on the
champion profile. At simulation time, reaction appeal is the niche weighted sum
of DNA dimensions rather than a single global dot product.

\subsection{Serving calibrated artifacts}
Community offsets are not part of the public \texttt{Profile} schema (they stay
in the private backend). During CMA-ES they attach as a runtime patch on the
profile object; on promotion we persist them in the artifact's metrics blob
alongside evaluation numbers. At every predict request, \texttt{active\_profile}
loads the champion payload \emph{and} merges stored community offsets (and any
author specific offsets from the audience table) onto the profile before
\texttt{analyze} runs. This guarantees that a promoted challenger changes live
behavior immediately, not only offline metrics.

\subsection{Learned reaction model}
\label{sec:reaction-mlp}
When enough OAuth authorized rows exist (impressions plus engagement breakdown),
a weekly job trains a small MLP to predict reaction logits from persona traits,
DNA, and niche membership. The best checkpoint is versioned like profile
artifacts. \texttt{engine\_service} loads the champion reaction model when
present and passes it into the cascade reaction step; otherwise the heuristic
path in Section~3.4 applies. Training and inference share the same feature
layout so promotion is a drop in replacement.

\subsection{Author audience automation}
After a user completes OAuth, a background job samples their recent posts and
public follower statistics (where available), estimates engagement rate
distributions, and upserts an \texttt{audience\_profiles} row. Resolution order
at predict time is: author specific profile, tenant generic profile, then global
champion birth priors only. This is how the same text scores differently for a
niche creator versus a mass media account without retraining the global profile.

\subsection{Causal experiments}
The experiments API allocates traffic between control and treatment arms---for
example two profile versions or two rewrite candidates from \texttt{/improve}.
Each arm logs a prediction id; maturation writes \texttt{rel\_amplification} per
arm. We report lift and confidence on the experiment record so product teams can
validate a challenger causally before full promotion.

\subsection{Drift monitoring and retrain schedule}
A daily monitor job recomputes holdout rank correlation, score saturation, label
coverage, and monotonicity checks on recent predictions. Breaches write
\texttt{monitoring\_events} rows (severity info/warning/critical) and can flag
an early retrain. Production crons run velocity snapshots every 30 minutes,
outcome maturation hourly, monitoring daily, and full profile plus reaction
retrain weekly. All jobs are idempotent module invocations behind API key auth
on a single Postgres store.

\paragraph{Promote.}
We split data by time: oldest \fill{75}\% train, newest \fill{25}\% test. A
challenger replaces the live champion only if it beats the champion on the
recent slice on several metrics at once, not only headline rank correlation. Every
version is versioned with full metrics for rollback.

\paragraph{Serve.}
The active champion profile and any stored offsets load on the next request with
no code deploy. Optional ensemble mode exposes uncertainty bands to the client.
The open source engine (\texttt{tweet\_virality\_simulator}) implements DNA,
niches, network, cascade, and scoring; the private service adds tracing,
Postgres, calibration, audience resolution, reaction model loading, and the
HTTP API described above.

\section{Deployment and open source split}

We separate concerns deliberately. The \textbf{open engine} is a pip installable
package with CLI commands \texttt{tvs x}, \texttt{compare}, \texttt{improve},
and \texttt{validate}. Anyone can score drafts locally with zero API keys. The
\textbf{private service} wraps the same \texttt{analyze} entry point, logs
predictions, ingests outcomes, runs the learning loop, and serves
\texttt{/predict} and \texttt{/improve} over HTTPS with API key authentication.

\paragraph{Data bootstrap.}
For bulk experiments we provide scripts that (i) score a CSV of tweets locally,
(ii) fetch a tweet by platform id and score it, and (iii) sample a public Kaggle
celebrity tweet dump, compute proxy labels, and ingest rows into the training
store via the predict endpoint. The same relative amplification definition is
used in bootstrap and in production maturation.

\paragraph{Infrastructure.}
The deployed stack is FastAPI + Uvicorn, Postgres (hosted on Supabase), and
cron workers on a single VM. Environment variables control outcome window hours,
minimum training sample count, optional compat LLM endpoints for DNA, and API
keys. Profile and reaction artifacts are rows in an \texttt{artifacts} table with
\texttt{champion} / \texttt{retired} status for rollback.

\section{Experimental Setup}
\label{sec:exp-setup}

\subsection{Data}
We use \fill{100{,}000} tweets drawn from \fill{N} public accounts spanning
\fill{date range}. \fill{Describe the source here, for example a public dataset
of tweets from well known accounts, and note any cleaning you did, such as
dropping pure retweets and non English posts.} For each author we compute the
baseline from their own history and label every tweet with its relative
amplification as defined above. We split the data by time into \fill{75} percent
for training and the most recent \fill{25} percent for testing, so that no future
information leaks into the fit.

\subsection{Baselines}
We compare the calibrated simulator against the following.
\begin{itemize}
  \item \textbf{Uncalibrated simulator.} The same engine with default
  parameters and no training. This isolates the value of the learning loop.
  \item \textbf{Content only heuristic.} A scorer that uses only the draft
  features without the cascade. This isolates the value of the simulation.
  \item \fill{Optional: a content based machine learning model trained on the
  same labels, for example a text classifier, as an external point of
  comparison.}
\end{itemize}

\subsection{Evaluation harness}
We evaluate with the same \texttt{evaluate\_profile} routine used in promotion:
Spearman rank correlation, log tail correlation, top decile recall, pairwise tier
accuracy, calibration MAE, and saturation. A holdout script replays a CSV of
\texttt{(tweet, rel\_amplification)} tuples through the active champion (with
offsets and niche weights loaded) and prints the metric bundle. Ablation rows
drop one component at a time---algorithmic channel, community offsets, niche
appeal, author audience profile, reaction MLP, sentence encoder---on the same
time split.

\subsection{Measures}
Because spread is heavy tailed, we report several measures rather than one.
\begin{itemize}
  \item \textbf{Rank correlation.} Spearman correlation between the predicted
  score and the real relative amplification. This is the headline number.
  \item \textbf{Tail rank correlation.} The same correlation computed in log
  space, which is more sensitive to the rare large hits.
  \item \textbf{Top group recall.} Of the tweets that really did best, how many
  the model placed in its own top group.
  \item \textbf{Pair accuracy.} For pairs of tweets in different tiers, how often
  the model ordered them correctly.
  \item \textbf{Calibration error.} How far the predicted level sits from the
  real level on average.
\end{itemize}

\section{Results}

\fill{Fill this section with your numbers. The tables below are ready to take
them.}

\begin{table}[t]
\centering
\caption{Ranking quality on the held out recent slice. Higher is better for all
columns except calibration error.}
\label{tab:main}
\begin{tabular}{lccccc}
\toprule
Model & Rank corr. & Tail rank & Top recall & Pair acc. & Calib. err. \\
\midrule
Content only        & \fill{ } & \fill{ } & \fill{ } & \fill{ } & \fill{ } \\
Uncalibrated sim.   & \fill{ } & \fill{ } & \fill{ } & \fill{ } & \fill{ } \\
Calibrated sim.\ (ours) & \fill{ } & \fill{ } & \fill{ } & \fill{ } & \fill{ } \\
\bottomrule
\end{tabular}
\end{table}

\noindent
Our main finding is that calibration moves the rank correlation from
\fill{Y.YY} to \fill{X.XX}, a gain of \fill{...}. On a 10{,}000--tweet sample
from the public Kaggle celebrity corpus (uncalibrated default profile, 40 Monte
Carlo runs per draft), the simulator already reaches Spearman $\rho = 0.718$
between score and relative amplification. The strongest per--tweet matches
(Table~\ref{tab:qual-examples}) are not always the highest--scoring drafts: a
\texttt{@jimmyfallon} primetime promo and a \texttt{@BarackObama} quote both sit
near the 83rd and 73rd percentiles of predicted and actual relative reach
respectively ($|\Delta| < 0.02$ percentile points), while a
\texttt{@jtimberlake} fan Q\&A bit ($4.5\times$ author median, score~44) is
ranked almost as accurately among the top decile of real hits.

\paragraph{Best predicted tweets.}
We rank every scored draft by how closely its percentile rank matches the
percentile of real relative amplification
($|\text{score pct.} - \text{rel.\ amp.\ pct.}|$; lower is better). Table~\ref{tab:qual-examples} lists the five lowest
error cases from that 10{,}000--tweet run.

\begin{table}[t]
\centering
\small
\caption{Qualitative examples where the simulator's ranking best matched real
relative reach. \emph{Rel.\ amp.} is engagement (likes $+ 2\times$ retweets)
divided by the author's own median; \emph{$|\Delta|$} is the absolute gap between
the draft's score percentile and its rel.\ amp.\ percentile within the sample.}
\label{tab:qual-examples}
\begin{tabular}{@{}p{0.11\textwidth}p{0.42\textwidth}cccr@{}}
\toprule
Author & Draft (excerpt) & Score & Rel.\ amp. & Likes & $|\Delta|$ \\
\midrule
\texttt{jimmyfallon} &
One week from tonight: The Best of Late Night 2 Hour Primetime Special\ldots &
24 & 2.35 & 2{,}556 & 0.01 \\
\texttt{BarackObama} &
``Justice is making sure that every young person knows that they are special
and their lives matter.'' ---President Obama &
15 & 1.70 & 2{,}542 & 0.01 \\
\texttt{jtimberlake} &
\#ThingsToNeverAskADJ ``You got any Backstreet Boys??'' Bam. &
44 & 4.52 & 1{,}627 & 0.02 \\
\texttt{ladygaga} &
The film documents the lives of young people affected by rape. Our song is
called ``Till It Happens To You.'' &
18 & 1.91 & 8{,}777 & 0.03 \\
\texttt{rihanna} &
Navy y'all did this i\$h!!! FIRST \& ONLY @RIAA certified with 100M song awards
\& \#ANTI is PLATINUM!! &
61 & 7.56 & 23{,}492 & 0.09 \\
\bottomrule
\end{tabular}
\end{table}

\noindent
\textbf{Primetime promo (\texttt{@jimmyfallon}).}
The draft announces a two--hour TV special and asks fans to set their DVRs.
Absolute likes (2{,}556) look large to a human reader, but relative to
\texttt{@jimmyfallon}'s own median the tweet is a solid but not explosive
$2.35\times$ hit. The simulator assigns score~24 (``Likely to fizzle'' in
absolute terms) yet places it at the 83rd percentile of both predicted and
actual relative reach---within 0.01 percentile points.

\noindent
\textbf{Fan Q\&A punchline (\texttt{@jtimberlake}).}
A short, humorous exchange (\#ThingsToNeverAskADJ) earned $4.5\times$ the
author's typical engagement, driven by 5{,}305 retweets. The simulator scores it
44 (``Limited spread'') and ranks it in the 94th percentile on both sides
($|\Delta| = 0.02$), capturing a sleeper hit that raw like counts alone would
not normalize across accounts.

\noindent
\textbf{Milestone announcement (\texttt{@rihanna}).}
Among the top--accuracy set, this is the clearest \emph{high}--reach case: rel.\
amp.\ $7.56\times$ with score~61 (``Has legs''), both near the 97th percentile.
The model correctly separates a genuine outlier from the long tail of
low--amplification replies that also receive high absolute scores on brand
accounts.

\begin{table}[t]
\centering
\caption{Effect of removing parts of the model. Each row drops one piece.}
\label{tab:ablation}
\begin{tabular}{lc}
\toprule
Setting & Rank corr. \\
\midrule
Full model                 & \fill{ } \\
\quad without the algorithmic channel & \fill{ } \\
\quad without per community offsets   & \fill{ } \\
\quad without niche appeal weights    & \fill{ } \\
\quad without author audience profile & \fill{ } \\
\quad without learned reaction MLP    & \fill{ } \\
\quad without sentence encoder (hash only) & \fill{ } \\
\quad without the cascade (features only) & \fill{ } \\
\bottomrule
\end{tabular}
\end{table}

\noindent
\fill{Describe the ablation. The point to make is which parts of the simulation
actually carry the signal.}

\paragraph{Speed.}
The simulator scores one draft in about \fill{Z} seconds with \fill{120} random
runs on \fill{your hardware}. \fill{If you have a number for a language model
based simulator, contrast it here to show the speed advantage.}

\section{Discussion}

The results suggest that a simple cascade, once its parameters are tied to real
outcomes, carries more signal about spread than the draft features alone. The
gain from the algorithmic channel in Table~\ref{tab:ablation} points to the same
thing that platform audits report, which is that out of network exposure now
drives a large share of reach.

Two design choices did most of the work. Training on relative amplification kept
the model honest about message quality, and the time based split kept us honest
about whether the model would help tomorrow rather than only explaining
yesterday. The staged stack in Table~\ref{tab:phases} matters as much as the
core cascade: community offsets and niche appeal capture heterogeneity that a
single global profile misses; author audience profiles personalize the crowd
without retraining the world model; the reaction MLP and sentence encoder add
capacity only where data and latency allow; monitoring and causal experiments
keep the loop honest after deployment.

\subsection{Limitations}
The text encoder is still weaker than a large language model, which may cap
subtle wording effects even with the sentence transformer option. We treat media
as a single flag rather than modeling image or video content directly. Public
engagement numbers remain a noisier stand in for impressions when OAuth is
unavailable. Our evaluation corpus comes from a particular set of accounts and
time window; platform ranking changes may shift absolute scores. The reaction
MLP and author audience jobs need sufficient per author data before they beat
heuristic defaults. Finally, the model predicts how far a message is likely to
spread, not whether it is true or worth spreading.

\section{Ethics}

A tool that estimates spread can be misused to optimize for outrage or to flood
feeds. Our aim is the opposite case, which is to give writers and analysts honest
feedback on a draft. We store outcome data behind access control, and the author
authorized path only reads metrics that the author has agreed to share. We do not
attempt to identify private individuals, and the model offers no advice about the
truth of a claim, only about its likely reach. We think the relative amplification
target is also a small safeguard, because it rewards saying something that lands
rather than simply having a large audience.

\section{Conclusion}

We set out to answer a question people ask before they post: will this travel?
Forecasting methods cannot, because they need a cascade that does not exist yet,
and content models tend to grade the account rather than the message. Our
simulator runs the crowd forward from a blank draft, scores how far the message
goes, and learns its parameters from real outcomes measured against each author's
own baseline. On \fill{100{,}000} tweets the trained version clearly improved
ranking over the untrained one while staying fast enough to run inside a web
request. The full continual learning stack---niche appeal, community offsets,
author audiences, reaction model, ensemble uncertainty, experiments, and
monitoring---closes the loop from draft to deployment without manual parameter
tuning. We release the open engine and describe the private service architecture
so others can reproduce the simulation core and extend the learning loop.

\section*{Reproducibility}

The open simulation engine is available at \fill{GitHub URL}. Install with
\texttt{pip install -e .} and run \texttt{tvs x "your draft"} locally. Bulk
scoring and Kaggle bootstrap scripts live under \texttt{examples/}. The private
service schema is in \texttt{sql/schema.sql} and \texttt{sql/migration\_v2.sql};
cron schedules are in \texttt{deploy/crontab.txt}. Before submission we will
publish \fill{(i)} a frozen champion profile artifact hash, \fill{(ii)} the
holdout evaluation CSV or access instructions for the \fill{100{,}000} tweet
corpus, and \fill{(iii)} a one command script that reproduces Table~\ref{tab:main}
and Table~\ref{tab:ablation} on the time split described in
Section~\ref{sec:exp-setup}.

% ---------------------------------------------------------------------------
\begin{thebibliography}{9}

\bibitem{zhao2015seismic}
Q.~Zhao, M.~Erdogdu, H.~He, A.~Rajaraman, and J.~Leskovec.
\newblock SEISMIC: A self exciting point process model for predicting tweet
popularity.
\newblock In \emph{Proc. KDD}, 2015.

\bibitem{cheng2014cascades}
J.~Cheng, L.~Adamic, P.~A.~Dow, J.~Kleinberg, and J.~Leskovec.
\newblock Can cascades be predicted?
\newblock In \emph{Proc. WWW}, 2014.

\bibitem{rameez2022viralbert}
R.~Rameez, H.~A.~Rahmani, and E.~Yilmaz.
\newblock ViralBERT: A user focused BERT based approach to virality prediction.
\newblock In \emph{Proc. UMAP}, 2022.

\bibitem{liu2025popsim}
Y.~Liu, W.~Liu, X.~Gu, A.~He, W.~Wang, and Y.~Zhang.
\newblock PopSim: Social network simulation for social media popularity
prediction.
\newblock \emph{arXiv preprint arXiv:2512.02533}, 2025.

\bibitem{park2023agents}
J.~S.~Park, J.~O'Brien, C.~J.~Cai, M.~R.~Morris, P.~Liang, and M.~S.~Bernstein.
\newblock Generative agents: Interactive simulacra of human behavior.
\newblock In \emph{Proc. UIST}, 2023.

\bibitem{weng2012competition}
L.~Weng, A.~Flammini, A.~Vespignani, and F.~Menczer.
\newblock Competition among memes in a world with limited attention.
\newblock \emph{Scientific Reports}, 2:335, 2012.

\bibitem{pezzoni2013retweet}
F.~Pezzoni, J.~An, A.~Passarella, J.~Crowcroft, and M.~Conti.
\newblock Why do I retweet it? An information propagation model for microblogs.
\newblock In \emph{Proc. SocInfo}, 2013.

\bibitem{huszar2022amplification}
F.~Husz\'ar, S.~I.~Ktena, C.~O'Brien, L.~Belli, A.~Schlaikjer, and M.~Hardt.
\newblock Algorithmic amplification of politics on Twitter.
\newblock \emph{PNAS}, 119(1), 2022.

\bibitem{hansen2016cma}
N.~Hansen.
\newblock The CMA evolution strategy: A tutorial.
\newblock \emph{arXiv preprint arXiv:1604.00772}, 2016.

\end{thebibliography}

\end{document}
